
\section*{Part I. Quantum optics \hfill 4 / 15}
\subsection*{1. Quadratures and optical phase space}
The course uses the dimensionless quadrature convention
\[
\boxed{\hat x=\hat a+\hat a^\dagger,\quad
\hat p=i(\hat a^\dagger-\hat a),\quad [\hat x,\hat p]=2i.}
\]
Consequently \((\Delta x)^2(\Delta p)^2\ge1\).
For \(\ket\alpha\):
\[
\langle x\rangle=2\operatorname{Re}\alpha,\quad
\langle p\rangle=2\operatorname{Im}\alpha,\quad
(\Delta x)^2=(\Delta p)^2=1.
\]
For phase angle \(\theta\),
\(\hat X_\theta=\hat a e^{-i\theta}+\hat a^\dagger e^{i\theta}\).
Balanced homodyne detection selects \(\theta\) via the local oscillator.
These formulas are convention-sensitive: some texts scale \(x,p\)
by \(\sqrt2\).

\subsection*{2. Squeezing changes noise, not the uncertainty bound}
For real \(r\), take
\(\hat S(r)=\exp[\frac r2(\hat a^2-\hat a^{\dagger2})]\).
The squeezed vacuum \(\hat S(r)\ket0\) has zero means but
\[
\boxed{(\Delta x)^2=e^{-2r},\qquad
(\Delta p)^2=e^{2r}.}
\]
One quadrature is quieter than vacuum and its conjugate noisier;
the product stays 1. Squeezing two appropriate modes and combining
them at a balanced beam splitter can produce entangled two-mode light.

\subsection*{3. Representing quantum states in phase space}
The Husimi function is always positive:
\(Q(\alpha)=\pi^{-1}\bra\alpha\rho\ket\alpha\).
The Glauber--Sudarshan \(P\) representation writes
\(\rho=\int P(\alpha)\ket\alpha\bra\alpha\,d^2\alpha\);
\(P\) may be singular or nonclassical. The Wigner function is a
real quasiprobability; negative regions provide a nonclassicality
witness. These are different representations of the same state.
For Gaussian states, a displacement vector
\(d=\langle R\rangle\), \(R=(x,p)^T\), and covariance
\(V_{jk}=\langle\{\Delta R_j,\Delta R_k\}\rangle/2\)
fully specify the state. A Gaussian optical circuit acts as
\(d'=Sd,\ V'=SVS^T\) (plus environmental noise for a channel).

\subsection*{4. Coherence and photon correlations}
\[
\boxed{g^{(2)}(0)=
\frac{\langle\hat a^{\dagger2}\hat a^2\rangle}{\bar n^2}
=\frac{\langle n(n-1)\rangle}{\bar n^2}.}
\]
Coherent light gives 1; single-mode thermal light gives 2;
a one-photon Fock state gives 0 (antibunching).
\begin{methodbox}
\textbf{Exam method:} identify the quadrature convention;
compute first and second moments or \(g^{(2)}\);
interpret the variance ellipse or photon statistics physically.
\end{methodbox}
\src{W1L3 pp.~8--27; W3L1 pp.~4--18; W3L2 pp.~3--21; W3L3 pp.~5--28.}
